\documentclass[11pt, a4paper]{report}

\usepackage[utf8]{inputenc}
\usepackage[T1]{fontenc}
\usepackage[french]{babel}

% Appel des packages dans le dossier preambule
\def\packagepath{../../../../../preambule}   % path du package princial

\usepackage{\packagepath/page_layout} % <--- Nouveau
\usepackage{\packagepath/config_info}
\usepackage{\packagepath/cadres_cours}
\usepackage{\packagepath/zones_reponse}
\usepackage{\packagepath/config_ds}
\usepackage{\packagepath/config_maths}

% --- PILOTAGE ---
%\def\visibleornot{visible}    % visible or invisible
\ifdefined\visibleornot\else\def\visibleornot{visible}\fi


\def\documentpath{./Sources_Latex}
\graphicspath{{./Images/}}


\def\ptsexoA{0}



\begin{document}

\def\level{Premiere}  % Classe à droite
\def\course{NSI}
\def\chaptername{Algèbre de Boole} % Titre au centre
\def\docname{RD02~--~TD 1}        % Identifiant à gauche

\pagestyle{DS_LP}

\NomPrenom{}

\bigskip

%%=============================================================




        \begin{exercice_cours}[Donner dans chaque cas la valeur du booleen \texttt{rep}:]

        \begin{tabular}{p{3.6cm}p{8.6cm}c}
    \hline
    Variables & Expression booléenne& Valeur de \verb|rep|\\
    \hline
        \verb|x = 3| & \verb|rep = (x**2 == 7)| & \reponseLigne[2.5cm]{\texttt{False}} \\
        \verb|x, y, z = 3, 4, 5| & \verb|rep = (x**2 + y**2 == z**2)| & \reponseLigne[2.5cm]{\texttt{False}} \\
        \verb|a, b = 3, -7| & \verb|rep = (a**3 > 50 and b**2 < 50)| & \reponseLigne[2.5cm]{\texttt{True}} \\
        \hline
        \end{tabular}
        \end{exercice_cours}

        \medskip

        \begin{exercice_cours}[Compléter la table de vérité suivante:]{}
            \begin{tabular}{cc|c|c|c|c|c|c|}
                \hline
                A & B & $\overline{A}$ & $\overline{B}$ & $\overline{A}\cdot B$ & $\overline{A} + B$ & $\overline{A} \cdot \overline{B}$ & $\overline{A} + \overline{B}$ \\
                \hline
                0 & 0 & \reponseLigne[1.5cm]{\texttt{1}} & \reponseLigne[1.5cm]{\texttt{1}} & \reponseLigne[1.5cm]{\texttt{0}} & \reponseLigne[1.5cm]{\texttt{1}} & \reponseLigne[1.5cm]{\texttt{1}} & \reponseLigne[1.5cm]{\texttt{1}} \\
                0 & 1 & \reponseLigne[1.5cm]{\texttt{1}} & \reponseLigne[1.5cm]{\texttt{0}} & \reponseLigne[1.5cm]{\texttt{1}} & \reponseLigne[1.5cm]{\texttt{1}} & \reponseLigne[1.5cm]{\texttt{0}} & \reponseLigne[1.5cm]{\texttt{1}} \\
                1 & 0 & \reponseLigne[1.5cm]{\texttt{0}} & \reponseLigne[1.5cm]{\texttt{1}} & \reponseLigne[1.5cm]{\texttt{0}} & \reponseLigne[1.5cm]{\texttt{0}} & \reponseLigne[1.5cm]{\texttt{0}} & \reponseLigne[1.5cm]{\texttt{1}} \\
                1 & 1 & \reponseLigne[1.5cm]{\texttt{0}} & \reponseLigne[1.5cm]{\texttt{0}} & \reponseLigne[1.5cm]{\texttt{0}} & \reponseLigne[1.5cm]{\texttt{1}} & \reponseLigne[1.5cm]{\texttt{0}} & \reponseLigne[1.5cm]{\texttt{0}} \\
                \hline
            \end{tabular}
        \end{exercice_cours}

        \medskip

        \begin{exercice_cours}[Créer la table de vérité de l'expression suivante: $S=(A \cdot B)+(A \cdot \overline{C})+(\overline{B}\cdot C)$]{}
            \begin{tabular}{ccc|P{10cm}|c|}
                \hline
                A & B & C & S \\
                \hline
                0 & 0 & 0 & &\reponseLigne[1.5cm]{\texttt{0}} \\
                0 & 0 & 1 & &\reponseLigne[1.5cm]{\texttt{1}} \\
                0 & 1 & 0 & &\reponseLigne[1.5cm]{\texttt{0}} \\
                0 & 1 & 1 & &\reponseLigne[1.5cm]{\texttt{0}} \\
                1 & 0 & 0 & &\reponseLigne[1.5cm]{\texttt{1}} \\
                1 & 0 & 1 & &\reponseLigne[1.5cm]{\texttt{1}} \\
                1 & 1 & 0 & &\reponseLigne[1.5cm]{\texttt{1}} \\
                1 & 1 & 1 & &\reponseLigne[1.5cm]{\texttt{1}} \\
                \hline
            \end{tabular}
            
        \end{exercice_cours}
    \end{document}
        \begin{minipage}{0.48\linewidth}
                \begin{exercice_cours}[Conversion Binaire -> Décimal]
        \begin{tabular}{p{3.6cm}c}
    \hline
    Nombre binaire & Nombre entier \\
    \hline
    \verb|10110110| & \reponseLigne[2.5cm]{\texttt{182}} \\
    \verb|01101101| & \reponseLigne[2.5cm]{\texttt{109}} \\
    \verb|00110101| & \reponseLigne[2.5cm]{\texttt{53}} \\
    \hline
        \end{tabular}
                \end{exercice_cours}
    \end{minipage}
%%=============================================================

\bigskip

\begin{minipage}{0.48\linewidth}

        \begin{exercice_cours}[Conversion Décimal -> Hexadécimal]

        \begin{tabular}{p{3.6cm}c}
    \hline
    Nombre entier & Nombre hexadécimal \\
    \hline
    \verb|42| & \reponseLigne[2.5cm]{\texttt{2A}} \\
    \verb|127| & \reponseLigne[2.5cm]{\texttt{7F}} \\
    \verb|64| & \reponseLigne[2.5cm]{\texttt{40}} \\
     \hline
        \end{tabular}
        \end{exercice_cours}
    \end{minipage}\hfill
    \begin{minipage}{0.48\linewidth}
                \begin{exercice_cours}[Conversion Hexadécimal -> Décimal]
        \begin{tabular}{p{3.6cm}c}
    \hline
    Nombre hexadécimal & Nombre entier \\
    \hline
    \verb|4A| & \reponseLigne[2.5cm]{\texttt{74}} \\
    \verb|7C| & \reponseLigne[2.5cm]{\texttt{124}} \\
    \verb|42| & \reponseLigne[2.5cm]{\texttt{66}} \\
    \hline
        \end{tabular}
                \end{exercice_cours}
    \end{minipage}

    %%%%%%%%%%

    
\begin{minipage}{0.48\linewidth}

        \begin{exercice_cours}[Conversion Binaire -> Hexadécimal]

        \begin{tabular}{p{3.6cm}c}
    \hline
    Nombre binaire & Nombre hexadécimal \\
    \hline
    \verb|10110101| & \reponseLigne[2.5cm]{\texttt{B5}} \\
    \verb|00011111| & \reponseLigne[2.5cm]{\texttt{1F}} \\
    \verb|10011010| & \reponseLigne[2.5cm]{\texttt{9A}} \\
     \hline
        \end{tabular}
        \end{exercice_cours}
    \end{minipage}\hfill
    \begin{minipage}{0.48\linewidth}
                \begin{exercice_cours}[Conversion Hexadécimal -> Binaire]
        \begin{tabular}{p{3.6cm}c}
    \hline
    Nombre hexadécimal & Nombre entier \\
    \hline
    \verb|FE| & \reponseLigne[2.5cm]{\texttt{254}} \\
    \verb|12| & \reponseLigne[2.5cm]{\texttt{18}} \\
    \verb|A3| & \reponseLigne[2.5cm]{\texttt{163}} \\
    \hline
        \end{tabular}
                \end{exercice_cours}
    \end{minipage}




\begin{exercice_cours}[Répondre aux questions suivantes:]

    \begin{itemize}
        \item Combien de valeurs différentes peut-on représenter avec 2 bits ? Avec 6 bits ? Avec 8 bits ? Avec 32 bits ? 
        \reponseLigne[0.8\linewidth]{\texttt{ respectivement 4, 64, 256 et 4294967296}}
        \item Combien de bits sont nécessaires pour représenter 1000 valeurs différentes?\hfill\reponseLigne[2cm]{\texttt{10 bits}}
        \item Quel est le plus grand entier naturel que l'on peut représenter avec 7 bits? \hfill\reponseLigne[2cm]{\texttt{127}}
    \end{itemize}

\end{exercice_cours}

\medskip

\begin{exercice_cours}[Effectuer les opérations binaires suivantes:]
    
    \begin{minipage}{0.48\linewidth}
       \begin{center}
       \texttt{101111} + \texttt{11011}  \reponseLigne[2.5cm]{=\texttt{1001010}} \\
       \end{center}
    \end{minipage}\hfill
    \begin{minipage}{0.48\linewidth}
        \begin{center}
       \texttt{110011101} + \texttt{11011}  \reponseLigne[2.5cm]{=\texttt{110101110}} \\
       \end{center}
    \end{minipage}

    \vspace{2cm}
\end{exercice_cours}

\medskip

\begin{exercice_cours}[Effectuer les opérations binaires suivantes:]
    
    \begin{minipage}{0.48\linewidth}
       \begin{center}
       \texttt{101111} $\times$ \texttt{1000}  \reponseLigne[2.5cm]{=\texttt{101111000}} \\
       \end{center}
    \end{minipage}\hfill
    \begin{minipage}{0.48\linewidth}
        \begin{center}
       \texttt{110011000} $\div $ \texttt{10}  \reponseLigne[2.5cm]{=\texttt{11001100}} \\
       \end{center}
    \end{minipage}

\end{exercice_cours}

\end{document}